\section{Examples of PEPS on the torus}\label{sec:peps_examples}

This section collects the two-dimensional examples of
\cite[Appendix~A, ``Two dimensions: PEPS'']{Cirac2021Matrix} whose states
can be computed on the discrete torus. Each example is a single tensor
$A^i_{\alpha\beta\gamma\delta}$ with one physical index $i$ and four virtual
indices, ordered top, right, down, left, placed at every site of a square
lattice. On the torus the four legs of the site $v=(x,y)$ are the edges from
$v$ to $(x,y+1)$ (top), from $v$ to $(x+1,y)$ (right), from $(x,y-1)$ to $v$
(down) and from $(x-1,y)$ to $v$ (left).

\begin{definition}[Torus PEPS of a four-leg tensor]%
    \label{def:peps_torus_site_tensor}
    \lean{TNLean.PEPS.torusDownEdge, TNLean.PEPS.torusLeftEdge,
        TNLean.PEPS.torusTopLeg, TNLean.PEPS.torusRightLeg,
        TNLean.PEPS.torusDownLeg, TNLean.PEPS.torusLeftLeg,
        TNLean.PEPS.torusLeftEdge_add_one, TNLean.PEPS.torusDownEdge_add_one,
        TNLean.PEPS.torusSiteTensor}
    \leanok
    \uses{def:peps_tensor, def:peps_torus_uniform_bond_dim}
    Let $a^i_{\alpha\beta\gamma\delta}$ be a tensor with physical dimension
    $d$ and bond dimension $D$. Its torus PEPS on the torus of width $w$ and
    height $h$ has bond dimension $D$ on every edge and, at every site $v$,
    the local tensor $a$ with its four virtual indices read on the top,
    right, down and left legs of $v$
    \cite[Appendix~A, ``Two dimensions: PEPS'']{Cirac2021Matrix}. The right
    leg of $v$ is the left leg of $v+(1,0)$, and the top leg of $v$ is the
    down leg of $v+(0,1)$.
\end{definition}

\begin{lemma}[Shift-invariant functions on the torus are constant]%
    \label{lem:peps_torus_shift_const}
    \lean{TNLean.PEPS.torusVertex_apply_eq_apply_zero_of_shift}
    \leanok
    \uses{def:peps_torus_geometry_translation}
    If a function $f$ on the sites of the torus satisfies $f(x+1,y)=f(x,y)$
    and $f(x,y+1)=f(x,y)$ for all $(x,y)$, then $f$ is constant.
\end{lemma}

\begin{proof}\leanok
    Induction on the representatives $0,\ldots,w-1$ of the first coordinate
    and $0,\ldots,h-1$ of the second.
\end{proof}

\begin{lemma}[Contraction of a torus PEPS]%
    \label{lem:peps_torus_site_tensor_contraction}
    \lean{TNLean.PEPS.stateCoeff_torusSiteTensor_eq_sum_edge,
        TNLean.PEPS.isHorizontalTorusEdge_torusRightEdge,
        TNLean.PEPS.isVerticalTorusEdge_torusUpEdge,
        TNLean.PEPS.torusRightEdge_injective, TNLean.PEPS.torusUpEdge_injective,
        TNLean.PEPS.torusRightEdge_ne_torusUpEdge, TNLean.PEPS.torusEdgeEquiv,
        TNLean.PEPS.stateCoeff_torusSiteTensor}
    \leanok
    \uses{def:peps_torus_site_tensor, def:peps_stateCoeff}
    Let $w,h\ge 3$. The right edges and the up edges of the sites are
    pairwise distinct and exhaust the edges of the torus, and the coefficient
    of the torus PEPS of $a$ is
    \begin{align}
        \Coeff(\sigma)=\sum_{h_\bullet,v_\bullet}\ \prod_{(x,y)}
        a^{\sigma(x,y)}_{v_{(x,y)}\,h_{(x,y)}\,v_{(x,y-1)}\,h_{(x-1,y)}},
        \notag
    \end{align}
    the sum running over all assignments of bond indices $h_{(x,y)}$ to the
    horizontal edges from $(x,y)$ to $(x+1,y)$ and $v_{(x,y)}$ to the
    vertical edges from $(x,y)$ to $(x,y+1)$.
\end{lemma}

\begin{proof}\leanok
    Every edge is horizontal or vertical, and every horizontal edge is the
    right edge of one of its endpoints and every vertical edge the up edge
    of one of its endpoints. For $w\ge 3$ the right edges of distinct sites
    are distinct: otherwise $p=q+(1,0)$ and $q=p+(1,0)$, so $2=0$ in the
    horizontal coordinate. The same holds vertically for $h\ge 3$. A
    horizontal edge is never vertical. Reindexing the sum over edge
    labellings along this bijection gives the formula.
\end{proof}

On a torus of width two the right and left neighbours of a site coincide and
the two edges joining them are one edge of the lattice graph, so the four
legs of a site are not four distinct bonds. The contraction formula, and the
cluster and CZX statements below, therefore require $w,h\ge 3$; the
source constructions are local and apply to every torus
\cite{gap:rmp_peps_examples_small_torus}.

\subsection{The GHZ state}

\begin{definition}[Two-dimensional GHZ tensor]%
    \label{def:peps_ghz_tensor}
    \lean{TNLean.PEPS.ghzSiteTensor, TNLean.PEPS.ghzPEPS}
    \leanok
    \uses{def:peps_torus_site_tensor}
    The GHZ tensor has $D=d$ and
    $A^i_{\alpha\beta\gamma\delta}=\delta_{i=\alpha=\beta=\gamma=\delta}$
    \cite[Appendix~A, ``The GHZ state'']{Cirac2021Matrix}.
\end{definition}

\begin{theorem}[The GHZ PEPS is the GHZ state]%
    \label{thm:peps_ghz_state}
    \lean{TNLean.PEPS.stateCoeff_ghzPEPS}
    \leanok
    \uses{def:peps_ghz_tensor, def:peps_stateCoeff}
    On every torus with $w,h\ge 2$ the GHZ PEPS generates the unnormalized
    GHZ state $\sum_{i=0}^{d-1}\ket{i,i,\ldots,i}$: its coefficient at
    $\sigma$ is $\sum_{i}\prod_{v}\delta_{\sigma(v),i}$
    \cite[Appendix~A, ``The GHZ state'']{Cirac2021Matrix}.
\end{theorem}

\begin{proof}\leanok
    \uses{lem:peps_torus_shift_const}
    A virtual configuration contributes only if every leg of every site
    carries the physical index of that site. The right leg of $v$ is the left
    leg of $v+(1,0)$ and the top leg of $v$ is the down leg of $v+(0,1)$, so
    $\sigma$ is invariant under both unit shifts and is constant by
    Lemma~\ref{lem:peps_torus_shift_const}. Every edge is then labelled by
    that constant, so exactly one virtual configuration contributes when
    $\sigma$ is constant and none otherwise.
\end{proof}

\subsection{The cluster state}

\begin{definition}[Two-dimensional cluster tensor]%
    \label{def:peps_cluster_tensor}
    \lean{TNLean.PEPS.clusterSiteTensor, TNLean.PEPS.clusterPEPS}
    \leanok
    \uses{def:peps_torus_site_tensor}
    The cluster tensor is $A=\ket{0}(00{+}{+}|+\ket{1}(11{-}{-}|$ with
    $(\pm|=[(0|\pm(1|]/\sqrt{2}$: the top and right legs copy the physical
    index $s$, and the down and left legs carry
    $(-1)^{sb}/\sqrt{2}$ at bond index $b$
    \cite[Appendix~A, ``The cluster state'']{Cirac2021Matrix}.
\end{definition}

\begin{theorem}[The cluster PEPS]%
    \label{thm:peps_cluster_state}
    \lean{TNLean.PEPS.stateCoeff_clusterPEPS}
    \leanok
    \uses{def:peps_cluster_tensor, def:peps_stateCoeff}
    On a torus with $w,h\ge 3$ and $N=wh$ sites the cluster PEPS has
    coefficient
    \begin{align}
        \Coeff(\sigma)=2^{-N}\prod_{v}(-1)^{\sigma_v\sigma_{v+(1,0)}
                +\sigma_v\sigma_{v+(0,1)}}.
        \notag
    \end{align}
\end{theorem}

\begin{proof}\leanok
    \uses{lem:peps_torus_site_tensor_contraction}
    By Lemma~\ref{lem:peps_torus_site_tensor_contraction} the top and right
    legs force every bond index to equal the physical index of the site
    below or to the left of it. The remaining factor at $v$ is
    $\tfrac12(-1)^{\sigma_v\sigma_{v-(0,1)}+\sigma_v\sigma_{v-(1,0)}}$;
    shifting the product over $v$ by one site in each direction gives the
    formula.
\end{proof}

Up to the factor $2^{-N/2}$ the coefficient in
Theorem~\ref{thm:peps_cluster_state} is that of the state obtained by
applying controlled-$Z$ gates on all nearest-neighbour bonds to
$\ket{+}^{\otimes N}$, the cluster state of
\cite[Appendix~A, ``The cluster state'']{Cirac2021Matrix}. The circuit and
this identification are not formalized.

\subsection{The CZX model}

In the CZX model of \cite{ChenLiuWen2011CZX} every site of the square lattice
carries four qubits, one at each corner, and the ground state on a closed
surface is the product over plaquettes of the states $\ket{0000}+\ket{1111}$
on the four qubits around each plaquette. Blocking the four qubits of a site,
\cite[Appendix~A, ``The CZX model'']{Cirac2021Matrix} writes this state as a
PEPS of bond dimension four.

\begin{definition}[CZX tensor]%
    \label{def:peps_czx_tensor}
    \lean{TNLean.PEPS.czxBond, TNLean.PEPS.czxQubits, TNLean.PEPS.czxTopLeft,
        TNLean.PEPS.czxTopRight, TNLean.PEPS.czxBottomRight,
        TNLean.PEPS.czxBottomLeft, TNLean.PEPS.czxSiteTensor,
        TNLean.PEPS.czxBondSwap, TNLean.PEPS.czxBondSwap_czxBond,
        TNLean.PEPS.czxPEPS}
    \leanok
    \uses{def:peps_torus_site_tensor}
    Label the four qubits of a site $\ket{ijkl}$, with $i$ at the top left,
    $j$ at the top right, $k$ at the bottom right and $l$ at the bottom left,
    and label the four basis vectors of a bond by pairs of qubits. The CZX
    tensor is
    \begin{align}
        \sum_{i,j,k,l=0}^{1}\ket{ijkl}\bigl((i,j),(j,k),(k,l),(l,i)\bigr|,
        \notag
    \end{align}
    so each leg carries the pair of corner qubits on its side, listed
    clockwise around the site
    \cite[Appendix~A, ``The CZX model'']{Cirac2021Matrix}. The CZX PEPS
    places this tensor at every site with its down and left labels read in
    reverse, so that each bond identifies the pair $(a,b)$ at one end with
    the pair $(b,a)$ at the other.
\end{definition}

The reversal corrects the printed formula. The review contracts every bond
with equal labels at its two ends
\cite[Section~``MPS and PEPS'']{Cirac2021Matrix}. Contracted in this way,
the printed tensor sets the top-right qubit of a site equal to the
bottom-left qubits of its right and upper neighbours, and the resulting state
is a product of GHZ states on diagonal loops that wind around the torus, not
the plaquette state. Two neighbouring sites list the qubits of their common
side in opposite orders, and the plaquette state is obtained by identifying
the pair $(a,b)$ at one end of a bond with $(b,a)$ at the other
\cite{gap:rmp_peps_czx_bond_orientation}.

\begin{theorem}[The CZX PEPS is a product of plaquette GHZ states]%
    \label{thm:peps_czx_plaquette_product}
    \lean{TNLean.PEPS.czx_bond_iff_plaquette, TNLean.PEPS.stateCoeff_czxPEPS}
    \leanok
    \uses{def:peps_czx_tensor, def:peps_stateCoeff}
    On a torus with $w,h\ge 3$ the CZX PEPS has coefficient
    \begin{align}
        \Coeff(\sigma)=\prod_{v}\delta\bigl[j_v=i_{v+(1,0)}=k_{v+(0,1)}
            =l_{v+(1,1)}\bigr],
        \notag
    \end{align}
    where $(i_v,j_v,k_v,l_v)$ are the four qubits of $\sigma(v)$ and the
    factor at $v$ is the coefficient of $\ket{0000}+\ket{1111}$ on the four
    qubits around the plaquette at the upper right corner of $v$
    \cite[Appendix~A, ``The CZX model'']{Cirac2021Matrix},
    \cite{ChenLiuWen2011CZX}.
\end{theorem}

\begin{proof}\leanok
    \uses{lem:peps_torus_site_tensor_contraction}
    By Lemma~\ref{lem:peps_torus_site_tensor_contraction} the top and right
    legs force the vertical bond above $v$ to carry $(i_v,j_v)$ and the
    horizontal bond to its right to carry $(j_v,k_v)$. The down and left legs
    then require $(i_{v-(0,1)},j_{v-(0,1)})=(l_v,k_v)$ and
    $(j_{v-(1,0)},k_{v-(1,0)})=(i_v,l_v)$ at every site, and these
    conditions, shifted by one site, are equivalent to the plaquette
    conditions.
\end{proof}

The review states that the CZX model lies in the non-trivial
symmetry-protected sector of its on-site $\mathbb{Z}_2$ symmetry
\cite[Appendix~A, ``The CZX model'']{Cirac2021Matrix}. This classification is
not formalized for the two-dimensional state. The boundary symmetry of the
model is a matrix product unitary, and the three-cocycle attached to a group
of matrix product unitaries is constructed in
Theorems~\ref{thm:mpug_anomaly_three_cocycle}
and~\ref{thm:mpug_anomaly_three_cocycle_gauge}.

\subsection{The toric code and quantum double models}

Let $G$ be a finite group. The quantum double model of $G$
\cite{Kitaev2003FaultTolerant} carries a spin with basis
$\{\ket{g}\}_{g\in G}$ on every edge of a square lattice; for
$G=\mathbb{Z}_2$ it is the toric code. The review gives two PEPS
representations \cite[Appendix~A, ``The Toric Code and quantum double
    models'']{Cirac2021Matrix}. Their symmetries are phrased through the map
$\mathcal{P}(A)=\sum A^i_{\alpha\beta\gamma\delta}\ket{i}\bra{\alpha\beta\gamma\delta}$
from the virtual to the physical system
\cite[Section~``MPS and PEPS'']{Schuch2010PEPS}.

\begin{definition}[$G$-injective and $G$-isometric tensors]%
    \label{def:peps_g_injective}
    \lean{TNLean.PEPS.IsGInjective, TNLean.PEPS.IsGIsometric,
        TNLean.PEPS.siteMap, TNLean.PEPS.siteMap_apply,
        TNLean.PEPS.apply_averageMap_of_forall_comp_eq,
        TNLean.PEPS.isGInjective_iff_exists_leftInverse,
        TNLean.PEPS.isGInjective_trivial_iff,
        TNLean.PEPS.siteMap_injective_iff}
    \leanok
    Let $g\mapsto U_g$ be a representation of a finite group $G$ on the
    virtual system of a tensor $A$, and let $\Pi_U$ be the projector onto the
    $U_g$-invariant subspace. The tensor is $G$-injective if
    $\mathcal{P}(A)U_g=\mathcal{P}(A)$ for every $g$ and $\mathcal{P}(A)$ is
    injective on the invariant subspace; equivalently, $\mathcal{P}(A)$ is
    invariant and has a left inverse $L$ with $L\,\mathcal{P}(A)=\Pi_U$
    \cite[Section~``Two dimensions: $G$-injective PEPS'']{Schuch2010PEPS}.
    It is $G$-isometric if moreover
    $\mathcal{P}(A)$ restricted to the invariant subspace is a multiple of an
    isometry, $\langle\mathcal{P}(A)x,\mathcal{P}(A)y\rangle
        =c\,\langle x,y\rangle$ for invariant $x,y$ and one constant $c>0$
    \cite[Section~``Isometric PEPS'']{Schuch2010PEPS}.
    For the trivial representation, $G$-injectivity is injectivity of
    $\mathcal{P}(A)$, that is, linear independence of the physical vectors
    $A^{\bullet}_{\alpha\beta\gamma\delta}$ indexed by the virtual
    configurations.
\end{definition}

The source requires $U_g$ to contain every irreducible representation for
$G$-injectivity and to be the left-regular representation for
$G$-isometry; here the representation is a parameter, and each instance
below uses a regular representation. The source's tensors are not
normalized, and the constant $c$ absorbs the normalization
\cite{gap:scp10_quantum_double_g_isometry}.

\begin{definition}[Dual quantum-double tensor]%
    \label{def:peps_quantum_double_dual}
    \lean{TNLean.PEPS.colorDiag, TNLean.PEPS.colorDiag_one,
        TNLean.PEPS.colorDiag_mul, TNLean.PEPS.quantumDoubleDualSpins,
        TNLean.PEPS.quantumDoubleDualTensor, TNLean.PEPS.colorShiftRep,
        TNLean.PEPS.colorShiftRep_apply,
        TNLean.PEPS.mem_invariants_colorShiftRep_iff}
    \leanok
    \uses{def:peps_g_injective}
    The virtual legs carry colors $h_1$ (left), $h_2$ (top), $h_3$ (right)
    and $h_4$ (down) in $G$, and the dual tensor is $1$ at the physical
    configuration
    $(h_2h_1^{-1},\,h_3h_2^{-1},\,h_3h_4^{-1},\,h_4h_1^{-1})$ of differences
    of adjacent colors and $0$ elsewhere \cite[Appendix~A, ``The Toric Code
        and quantum double models'']{Cirac2021Matrix}. The group acts on the
    virtual system by shifting all four colors,
    $(U_gx)(h_1,h_2,h_3,h_4)=x(h_1g,h_2g,h_3g,h_4g)$, the right-regular
    representation on each leg, which the relabeling $h\mapsto h^{-1}$
    turns into the left-regular one
    \cite{gap:scp10_quantum_double_g_isometry}.
\end{definition}

\begin{theorem}[The dual tensor is $G$-isometric]%
    \label{thm:peps_quantum_double_dual_g_isometric}
    \lean{TNLean.PEPS.quantumDoubleDualSpins_mul_colorDiag,
        TNLean.PEPS.quantumDoubleDualSpins_eq_iff,
        TNLean.PEPS.siteMap_quantumDoubleDualTensor_apply,
        TNLean.PEPS.siteMap_quantumDoubleDualTensor_apply_spins,
        TNLean.PEPS.siteMap_quantumDoubleDualTensor_apply_spins_of_mem_invariants,
        TNLean.PEPS.isGInjective_quantumDoubleDualTensor,
        TNLean.PEPS.isGIsometric_quantumDoubleDualTensor,
        TNLean.PEPS.ToricCodeGroup,
        TNLean.PEPS.isGIsometric_toricCodeDualTensor}
    \leanok
    \uses{def:peps_quantum_double_dual}
    For every finite group $G$ the dual quantum-double tensor is
    $G$-injective and $G$-isometric, with constant $c=|G|$, for the
    color-shift representation
    \cite[Appendix~A, ``The Toric Code and quantum double
        models'']{Cirac2021Matrix},
    \cite[Section~``Examples'', ``The double models'']{Schuch2010PEPS}.
    In particular the toric code tensor, $G=\mathbb{Z}_2$, is
    $\mathbb{Z}_2$-isometric.
\end{theorem}

\begin{proof}\leanok
    Two colorings have the same differences exactly when they differ by one
    global shift $h_i\mapsto h_ig$: the shift of the left color determines
    $g$, and the other three differences then fix the remaining colors.
    Hence $(\mathcal{P}x)(\text{differences of }h)=\sum_{g}x(hg)$. Invariance
    follows by reindexing the sum over colorings by a shift. For an
    invariant $x$ the sum equals $|G|\,x(h)$, so $\mathcal{P}x=0$ forces
    $x=0$, and
    $\langle\mathcal{P}x,\mathcal{P}y\rangle
        =\sum_h\overline{(\mathcal{P}x)(\text{differences of }h)}\,y(h)
        =|G|\langle x,y\rangle$.
\end{proof}

The review adds that the equal-weight superposition of plaquette colorings
is the equal-weight superposition of the configurations obeying Gauss' law.
This statement about the contracted network is not formalized; on a torus
the differences of plaquette colorings are the Gauss-law configurations of
trivial holonomy \cite{gap:scp10_quantum_double_g_isometry}.

\begin{definition}[Primal quantum-double tensor]%
    \label{def:peps_quantum_double_primal}
    \lean{TNLean.PEPS.quantumDoublePrimalLabels,
        TNLean.PEPS.quantumDoublePrimalTensor}
    \leanok
    A tensor carries the four edge spins $g_1,g_2,g_3,g_4$ of a plaquette (top
    left, top right, bottom right, bottom left), and is $1$ when its virtual
    labels are $g_1g_2$ (top), $g_2^{-1}g_3$ (right), $g_4g_3$ (down) and
    $g_1g_4^{-1}$ (left), and $0$ otherwise \cite[Appendix~A, ``The Toric Code
        and quantum double models'']{Cirac2021Matrix}.
\end{definition}

\begin{theorem}[The primal virtual labels fuse to the identity]%
    \label{thm:peps_quantum_double_primal_fusion}
    \lean{TNLean.PEPS.quantumDoublePrimalTensor_ne_zero_mul_eq_one,
        TNLean.PEPS.quantumDoublePrimalTensor_ne_zero_map_mul_eq_one,
        TNLean.PEPS.quantumDoublePrimalTensor_smul_map_mul,
        TNLean.PEPS.character_mul_quantumDoublePrimalTensor}
    \leanok
    \uses{def:peps_quantum_double_primal}
    On the support of the primal tensor the virtual labels satisfy
    $t\,r\,b^{-1}l^{-1}=1$. Hence $\pi(t)\pi(r)\pi(b)^{-1}\pi(l)^{-1}=1$ for
    every representation $\pi$ of $G$, in particular every irreducible one,
    so that each entry of the tensor times the operator
    $\pi(t)\pi(r)\pi(b)^{-1}\pi(l)^{-1}$, acting by $\pi$ on the top and
    right legs and by $\pi^{-1}$ on the down and left legs, equals the same
    entry times the identity. For a one-dimensional representation $\chi$
    the tensor is invariant under this virtual action
    \cite[Appendix~A, ``The Toric Code and quantum double
        models'']{Cirac2021Matrix}. This is a statement about one tensor; the
    resulting symmetry of the contracted network is not formalized
    \cite{gap:scp10_quantum_double_g_isometry}.
\end{theorem}

\begin{proof}\leanok
    $(g_1g_2)(g_2^{-1}g_3)(g_4g_3)^{-1}(g_1g_4^{-1})^{-1}
        =g_1g_3g_3^{-1}g_4^{-1}g_4g_1^{-1}=1$.
\end{proof}

\subsection{Examples not formalized}

The remaining two-dimensional examples of
\cite[Appendix~A, ``Two dimensions: PEPS'']{Cirac2021Matrix} are recorded
here without formal statements.
\begin{itemize}
    \item The two-dimensional AKLT state places singlets on the edges of the
          lattice and projects onto the symmetric subspace at each site; on the
          square lattice the tensor is $\Pi_{\mathrm{sym}}(\mathbb{1}\otimes
              \mathbb{1}\otimes Y\otimes Y)$ with $Y$ the singlet matrix.
    \item The nearest-neighbour resonating valence bond state is the
          superposition of all nearest-neighbour singlet coverings, a PEPS of bond
          dimension three built from a projector onto configurations with exactly
          one non-vacuum leg together with singlet matrices on the edges.
    \item String-net models give PEPS tensors built from the $F$-symbols of a
          bimodule category.
    \item A classical model with nearest-neighbour interactions at inverse
          temperature $\beta$ gives the PEPS
          $\bigl[\sum_i\ket{i}(i\,i\,i\,i|\bigr](\mathbb{1}\otimes\mathbb{1}
              \otimes M\otimes M)$ with $M=\sum_{i,j}e^{-\beta h(i,j)/2}|i)(j|$,
          whose diagonal expectation values reproduce those of the Gibbs state.
\end{itemize}
